\paragraph{All full types above a profile maximum.}
For a greatest profile completion, inherited singleton profiles completely
determine which sound source boxes can occur at one marked leaf. This also
gives the exact distinguished-source count, refining the interval bound in
\cref{core:thm:profile-source-interval}.

Fix finite $K\geq3$ and a greatest profile completion $C$ on a finite nonempty passive support
$W$, write $f_a=\Phi_a(C)$, and let $F\in\mathcal I(C)$ have distinguished
predicate $P_F$. A normalized full type $B=(B_r)_r$ consists of antitone
factors $B_r:2^W\to\{0,1\}$ with $B_r(\varnothing)=1$. Put
\[
 \psi_r=\neg B_r,\qquad R=\{r:\psi_r\ne0\},\qquad
 P_B(E)=\bigwedge_rB_r(E_r).
\]
Call the box sound for $F$ when $P_B\leq P_F$ on proper colorings.
On a primitive one-$X$ spine, $R$ records exactly the owners occurring on
the marked path: every nonempty passive helper has a positive top profile
at its gate's owner.

\begin{theorem}[Exact interval occurrence types]
\label{core:thm:profile-interval-types}
The full type $B$ is actual for $F$ if and only if its box is sound and
the following conditions hold:
\begin{enumerate}
\item $f_r\leq\psi_r$ for each $r\in R$;
\item for each $r\in R$, every coalition $A$ with $\psi_r(A)=1$ and
  every choice of coalitions $A_a$ with $f_a(A_a)=1$ for the inactive
  owners $a\notin R$ satisfy
  \[
   A\cap\bigcap_{a\notin R}A_a\ne\varnothing.
  \]
\end{enumerate}
When these hold, there is a primitive spine $S\subseteq F$ containing
$X$ once, with full type exactly $B$ and at most $2|R|$ owned nodes on
its marked path. This bound counts path attachments, not gates inside
the passive helpers.
\end{theorem}
\begin{proof}
For necessity linearize an actual occurrence by
\cref{core:cor:interval-linear-cover}, obtaining $S\subseteq F$ with
the same type. \Cref{core:lem:linear-profile-inheritance} gives
condition~(1) and makes every helper inherit the inactive profiles.
For $r\in R$, the full factor formula gives
\[
 \psi_r=\bigvee_{\text{$r$-owned path gates}}\Phi_r(H).
\]
Choose an $r$-helper witnessing $\psi_r(A)=1$. Its inherited profiles
at all inactive owners and its top profiles at the other active owners
give condition~(2) by profile admissibility.

Soundness follows from $S\subseteq F$ at source probes.

Conversely, condition~(2) makes the profile tuple
\[
 \Phi_r(H_r)=\psi_r,\qquad
 \Phi_a(H_r)=f_a\ (a\notin R),\qquad
 \Phi_s(H_r)(A)=[A=W]\ (s\in R\setminus\{r\})
\]
admissible. Profile synthesis supplies a passive core helper $H_r$ for
each $r\in R$. Start at $X$, attach $H_r$ once at owner $r$ for every
active owner, and then attach $C$ once at every active owner, using any
fixed orders in the two phases. The first phase gives exactly $B$.
The restoring phase adds factors $\neg f_r$, which are redundant by
condition~(1).

To check the bound, if $C=\{s\}$ with $s\in R$, its $s$-stage in
the restoring phase absorbs the prefix; subsequent copies of $C$ preserve
that singleton. If $C=\{a\}$ with $a\notin R$, every helper is
$\{a\}$ by inheritance, so the entire spine stays in $X\cup\{a\}$.
If $C$ is ambiguous, it is the full passive union, containing every
helper's value, so the spine stays in $C\cup X$ as well.
Thus $C+S$ lies in the interval and has distinguished predicate $P_B$.
Soundness gives $C+S\subseteq F$.

Restore $F$ by derived union with $S$. This preserves the marked type:
at a pure-$r$ source probe allowed by $B_r$, both operands contain the
queried foreign winner and hence neither is the $r$ singleton. No
subword of the union implementation can create that singleton, so its
added $r$-filters are redundant. The marked type is therefore actual.
If $R=\varnothing$, use $S=X$ and the same restoration argument.
\end{proof}

The profile parameters are independent once the active set is fixed.
For $J\subseteq I$, put
\[
 h_J(A)=1\quad\Longleftrightarrow\quad
 A\cap\bigcap_{a\in J}A_a\ne\varnothing
 \text{ for every choice with }f_a(A_a)=1.
\]
The empty intersection is $W$. These are normalized monotone profiles.
Before imposing soundness, the admissible tuples with active set $R$
are exactly the independent choices of normalized monotone profiles
\[
 f_r\leq\psi_r\leq h_{I\setminus R}\quad(r\in R),
 \qquad \psi_a=0\quad(a\notin R).
\]
Each interval is nonempty by admissibility of the baseline profiles.

For a predicate $P$, let $b_C(P)$ be the least number of boxes $P_B\leq P$
whose union is $P$ and whose factors satisfy (1)--(2). The empty cover has size zero.
For $P=P_F$ a finite cover exists by the theorem, using the types of any
realizer.

\begin{corollary}[Exact relative source multiplicity]
\label{core:cor:profile-interval-copy-number}
The minimum number of $X$-occurrences in a realizer of $F$ with union
named is $b_C(P_F)$.
\end{corollary}
\begin{proof}
Every occurrence contributes one sound box, and these boxes cover $P_F$.
The necessity argument above applies also with union named: extend marked
linearization by skipping a union node on the marked path. Union adds no
owner factor, and enlarging both the word and its residual preserves the
linearization bound. Thus all these boxes satisfy (1)--(2).

Conversely take a minimum cover, construct its primitive one-$X$ spines,
and union them with $C$. The bounds and source reconstruction give exactly
$F$, with one occurrence per box. If $P_F=0$, use $C$ alone.
\end{proof}

\begin{corollary}[A criterion for a greatest full type]
\label{core:cor:profile-interval-greatest-type}
Assume $P_F\ne0$, and recover the pure factors
$D_r(E)=P_F(\text{only }E\text{ colored }r)$.
There is a greatest actual full type if and only if $(D_r)_r$ is sound
and satisfies (1)--(2). Soundness here is equivalent to separability of
the distinguished source family.
\end{corollary}
\begin{proof}
Pure probes bound every actual factor by $D_r$. In any realizer, the
pointwise union of the marked $r$-factors is $D_r$ itself. A greatest
actual type would therefore equal $D$, and the theorem characterizes
exactly when $D$ is actual. Antitonicity gives
$P_F\leq\bigwedge_rD_r$; soundness supplies the reverse inequality.
\end{proof}

\paragraph{General baselines need simultaneous budgets.}
The profile-maximum hypothesis cannot be replaced by separate complete
bounded-profile tests, even for one active owner.

\begin{corollary}[Pointwise helper budgets do not determine actual types]
\label{core:cor:pointwise-helper-obstruction}
Take $C$ from \eqref{core:eq:profile-gluing-words}. In the core interval
operation $F(T,w)=C(w)+T$, the full type
\[
 B_c(E)=[\{x,y\}\nsubseteq E],\qquad B_r(E)=1\quad(r\ne c)
\]
is sound and satisfies active-owner profile containment and admissibility
with the inherited inactive-owner profiles. At every passive input, those
exact helper profiles have a realizer bounded by $C$ there. The type is
nevertheless not actual.
\end{corollary}
\begin{proof}
The inherited profile family is exactly that of
\cref{core:prop:profile-bound-gluing}, so the pointwise
claims follow from \cref{core:prop:profile-bound-gluing}; soundness follows
because $F$ retains $T$. If the type were actual, marked linearization
(\cref{core:cor:interval-linear-cover}) would give a primitive one-$T$
spine below $C+T$ with that exact type. Its path has only owner $c$:
every passive helper has a positive top profile at its gate's owner.
Associativity therefore writes it as $T\star_c H$, with
$\Phi_c(H)=xy$. Its neutral residual is $H\subseteq C$, contradicting
\cref{core:prop:profile-bound-gluing}.

The obstruction also applies if helper bounds are imposed only where
$\{x,y\}\nsubseteq E_c(w)$. On the omitted face $x=y=\{c\}$, every
helper with $c$-profile $xy$ equals $\{c\}$ and
$C=\{c\}\cup z$, so the remaining bounds hold automatically.
\end{proof}
